%------------------------------------------------------------------------------%
% Luis A. D'Afonseca
%
%------------------------------------------------------------------------------%
\documentclass[a4paper,12pt,fleqn]{article}

\usepackage{style_avaliacao}

\ShowAnswers

%------------------------------------------------------------------------------%
\begin{document}

\makeHeader{Integrais e Séries}{Prova 4 -- Turma 8}{14/02/2025}


% 01
%------------------------------------------------------------------------------%
\exercise{25}
Escreva o polinômio de Taylor de quarto grau, centrado no zero,
da função cosseno hiperbólico
\(
  \cosh(x)=\frac{e^x+e^{-x}}{2}
\)

\begin{answer}
  O polinômio de Taylor de quarto grau, centrado em zero, é
  \[
    P_4(x)
    = \cosh(0)
    + \cosh'(0) x
    + \frac{\cosh''(0)}{2}x^2
    + \frac{\cosh^{(3)}(0)}{3!}x^3
    + \frac{\cosh^{(4)}(0)}{4!}x^4
  \]
  Calculando as derivadas
  \begin{align*}
    \cosh^{(1)}(x)
    & = \frac{1}{2}\frac{d}{dx}\left(e^x+e^{-x}\right)
      = \frac{1}{2}\left(e^x-e^{-x}\right) \\[2mm]
    \cosh^{(2)}(x)
    & = \frac{1}{2}\frac{d}{dx}\left(e^x-e^{-x}\right)
      = \frac{1}{2}\left(e^x+e^{-x}\right) \\[2mm]
    \cosh^{(3)}(x)
    & = \frac{1}{2}\frac{d}{dx}\left(e^x+e^{-x}\right)
      = \frac{1}{2}\left(e^x-e^{-x}\right) \\[2mm]
    \cosh^{(4)}(x)
    & = \frac{1}{2}\frac{d}{dx}\left(e^x-e^{-x}\right)
      = \frac{1}{2}\left(e^x+e^{-x}\right)
  \end{align*}
  Avaliando as derivadas em zero
  \begin{align*}
    \cosh^{(0)}(0) & = \frac{e^0+e^{-0}}{2} = \frac{1+1}{2} = 1  \\[2mm]
    \cosh^{(1)}(0) & = \frac{e^0-e^{-0}}{2} = \frac{1-1}{2} = 0  \\[2mm]
    \cosh^{(2)}(0) & = \frac{e^0+e^{-0}}{2} = \frac{1+1}{2} = 1  \\[2mm]
    \cosh^{(3)}(0) & = \frac{e^0-e^{-0}}{2} = \frac{1-1}{2} = 0  \\[2mm]
    \cosh^{(4)}(0) & = \frac{e^0+e^{-0}}{2} = \frac{1+1}{2} = 1
  \end{align*}
  Portanto
  \[
    P_4(x)
    = 1
    + \frac{x^2}{2}
    + \frac{x^4}{4!}
  \]
\end{answer}

% 02
%------------------------------------------------------------------------------%
\exercise{25}
Considerando a série de potências
\(
  \sum_{n=0}^\infty \frac{x^{2n+1}}{(2n+1)n!}
\)
\begin{enumerate}[label=\alph*)]
  \item\ [15] determine seu raio de convergência
  \item\ [10] calcule sua derivada
\end{enumerate}

\begin{answer}
\begin{enumerate}[label=\alph*)]
\item
  Usando o teste da razão
  \begin{align*}
    a_n & = \frac{x^{2n+1}}{(2n+1)n!} &
    \abs{a_n} & = \frac{\abs{x}^{2n+1}}{(2n+1)n!} \\
    a_{n+1} & = \frac{x^{2n+3}}{(2n+3)(n+1)!} &
    \abs{a_n} & = \frac{\abs{x}^{2n+3}}{(2n+3)(n+1)!}
  \end{align*}
  \begin{align*}
    \rho
    & = \lim_{n\to\infty} \frac{\abs{a_{n+1}}}{\abs{a_n}} \\[2mm]
    & = \lim_{n\to\infty} \frac{\abs{x}^{2n+3}}{(2n+3)(n+1)!}\;
                          \frac{(2n+1)n!}{\abs{x}^{2n+1}}\\[2mm]
    & = \lim_{n\to\infty} \frac{2n+1}{(2n+3)(n+1)} \abs{x}^2 \\[2mm]
    & =  \abs{x}^2 \lim_{n\to\infty} \frac{2n+1}{2n^2+5n+3} \\[2mm]
    & =  \abs{x}^2 \lim_{n\to\infty} \frac{\sfrac{2}{n}+\sfrac{1}{n^2}}{2+\sfrac{5}{n}+\sfrac{3}{n^2}} \\[2mm]
    & = 0
  \end{align*}
  Portanto o raio de convergência é infinito \(R=\infty\)

\item
  \begin{align*}
    f(x)
    & = \frac{d}{dx} \sum_{n=0}^\infty \frac{x^{2n+1}}{(2n+1)n!} \\[2mm]
    & = \sum_{n=0}^\infty \frac{1}{(2n+1)n!} \frac{d}{dx} x^{2n+1} \\[2mm]
    & = \sum_{n=0}^\infty \frac{1}{(2n+1)n!} (2n+1)x^{2n} \\[2mm]
    & = \sum_{n=0}^\infty \frac{x^{2n}}{n!}
  \end{align*}
  Curiosidade
  \[
    \sum_{n=0}^\infty \frac{x^{2n}}{n!}
    = \sum_{n=0}^\infty \frac{\left(x^2\right)^n}{n!}
    = e^{2x}
  \]
\end{enumerate}
\end{answer}

% 03
%------------------------------------------------------------------------------%
\exercise{25}
Seja \( f(x) = \ln(x+1)\)
\begin{enumerate}[label=\alph*)]
  \item\ [10] Encontre um padrão para \(f^{(n)}(x)\)
  \item\ [15] Usando o padrão do item anterior, encontre a
              Serie de Maclaurin de $f(x)$
\end{enumerate}

\begin{answer}
  \begin{enumerate}[label=\alph*)]
    \item Calculando as primeiras derivadas
    \begin{align*}
      f^{(0)}(x) & = \ln(x+1) \\[2mm]
      f^{(1)}(x) & = \frac{\,d\;}{\,dx\;}\ln(x+1) = \frac{1}{x+1}(x+1)' = (x+1)^{-1} \\[2mm]
      f^{(2)}(x) & = \frac{d^2}{dx^2}    \ln(x+1) =    -                  (x+1)^{-2} \\[2mm]
      f^{(3)}(x) & = \frac{d^3}{dx^3}    \ln(x+1) = \hpm 2              \,(x+1)^{-3} \\[2mm]
      f^{(4)}(x) & = \frac{d^4}{dx^4}    \ln(x+1) =    - 2\cdot 3       \,(x+1)^{-4} \\[2mm]
      f^{(5)}(x) & = \frac{d^5}{dx^5}    \ln(x+1) = \hpm 2\cdot 3\cdot 4\,(x+1)^{-5}
    \end{align*}
    Vemos que, para $\bm{n\geq1}$,
    \[
      f^{(n)}(x)
      = \frac{d^n}{dx^n} \ln(x+1)
      = (-1)^{n+1} (n-1)! (x+1)^{-n}
      = \frac{(-1)^{n+1} (n-1)!}{(x+1)^n}
    \]
    \item O primeiro coeficiente, $n=0$, da Série de Maclaurin é
      \[
        c_0 = \frac{\ln(0+1)}{0!} = \frac{0}{1} = 0
      \]
      portanto, a série começa com $n=1$.
      Os demais coeficientes, $n\geq 1$, são
      \[
        c_n
        = \frac{f^{(n)}(0)}{n!}
        = \frac{(-1)^{n+1} (n-1)!}{(0+1)^nn!}
        = \frac{(-1)^{n+1}}{n}
      \]
      Montando a série temos
      \[
        T(x)
        = \sum_{n=0}^\infty c_n x^n
        = \sum_{n=1}^\infty \frac{(-1)^{n+1}}{n} x^n
      \]
  \end{enumerate}
\end{answer}

% 04
%------------------------------------------------------------------------------%
\exercise{25}
Use o polinômio de Taylor de segundo grau, centrado em 1, para aproximar o
valor de \(\sqrt{2}\), estime o erro cometido.
Não é necessário encontrar a representação decimal dos valores numéricos.
Dica: O valor máximo vai acontecer em um dos pontos extremos do intervalo.

\begin{answer}
  Polinômio de Taylor
  \[
    P_2(x) = f^{(0)}(1) +  f^{(1)}(1)(x-1) +  \frac{f^{(2)}(1)}{2} (x-1)^2
  \]
  Calculando as derivadas
  \begin{align*}
    f^{(0)}(x)
    & = \sqrt{x}
      = x^{\sfrac{1}{2}} \\[2mm]
    f^{(1)}(x)
    & = \frac{d}{dx} x^{\sfrac{1}{2}}
      = \frac{1}{2} x^{\sfrac{1}{2}-1}
      = \frac{x^{-\sfrac{1}{2}}}{2} \\[2mm]
    f^{(2)}(x)
    & = \frac{d}{dx} \frac{x^{-\sfrac{1}{2}}}{2}
      = \frac{1}{2}\; \frac{-1}{2}\;  x^{-\sfrac{1}{2}-1}
      = \frac{-x^{-\sfrac{3}{2}}}{4}
  \end{align*}
  Avaliando para $x=1$
  \begin{align*}
    f^{(0)}(1) & = \sqrt{1} = 1 \\[2mm]
    f^{(1)}(1) & = \frac{x^{-\sfrac{1}{2}}}{2} \evalat{x=1}{} = \frac{1}{2}\\[2mm]
    f^{(2)}(1) & = \frac{-x^{-\sfrac{3}{2}}}{4} \evalat{x=1}{} = \frac{-1}{4}
  \end{align*}
  Então
  \begin{align*}
    P_2(x)
    & = f(1) + f'(1)(x-1) + \frac{f''(1)}{2!}(x-1)^2 \\
    & = 1 + \frac{x-1}{2} -\frac{(x-1)^2}{8}
  \end{align*}
  Avaliando em 2
  \begin{align*}
    P_2(2)
    & = 1 + \frac{2-1}{2} - \frac{(2-1)^2}{8} \\
    & = 1 + \frac{1}{2}   - \frac{1}{8} \\
    & = \frac{8+4-1}{8} \\
    & = \frac{11}{8}
  \end{align*}


  Pelo \textbf{Teorema de Taylor}, sabemos que o erro cometido ao aproximarmos
  $\sqrt{x}$ pelo seu polinômio de Taylor centrado em 1 de grau 2 é
  \[
    R_2(x) = \frac{f^{(3)}(c)}{3!} \, (x-1)^3
  \]
  para algum $c$ entre zero e $x$ e
  \[
    f^{(3)}(x)
    = \frac{d}{dx} \frac{-x^{-\sfrac{3}{2}}}{4}
    = \frac{-1}{4} \; \frac{-3}{2} \; x^{-\sfrac{3}{2}-1}
    = \frac{3}{8}\; x^{-\sfrac{5}{2}}
    = \frac{3}{8\sqrt{x^5}}
  \]
  Entre 1 e 2 o maior valor de \(\abs{f^{(3)}(x)}\) acontece em $x=1$, portanto
  \[
    \abs{f^{(3)}(x)}
    \leq \frac{3}{8\sqrt{x^5}} \evalat{x=1}{}
    = \frac{3}{8}
  \]
  \[
    \abs{R_2(2)}
    \leq \abs{\frac{f^{(3)}(c)}{3!} \, (2-1)^3}
    = \frac{\sfrac{3}{8}}{3!}
    = \frac{3}{8\cdot 3\cdot 2\cdot 1}
    = \frac{1}{16}
  \]
\end{answer}

%------------------------------------------------------------------------------%
\end{document}
%------------------------------------------------------------------------------%
